% =============================================================================
\section{Software Architecture \& Component Design}
\label{sec:arch}
% =============================================================================

\subsection{Architectural Paradigm: Layered Abstraction and Dependency Inversion}

The firmware is structured strictly according to the \textbf{Dependency Inversion Principle (DIP)}
from modern safety-critical software engineering. Rather than having high-level business logic
(the FSM) directly read peripheral registers or invoke hardware-specific drivers, both high-level
and low-level components depend on abstract interfaces defined in \code{hal_interfaces.h}.

Figure~\ref{fig:arch} illustrates the four-tier architectural hierarchy:
\begin{enumerate}
  \item \textbf{Tier 1: Physical Hardware \& Microcontroller Peripherals:} Raw GPIO registers, SysTick
        counter, SPI4 transmit engine, and USART1 UART peripheral.
  \item \textbf{Tier 2: Hardware Abstraction Layer (HAL) Drivers:} Platform-independent algorithms
        for 30\,ms switch debounce integration, non-blocking LED waveform
        synthesis, multi-gesture decoding, and pulse-train coin recognition.
  \item \textbf{Tier 3: Core Finite State Machine (FSM):} The deterministic state engine,
        maintaining zero dependencies on microcontroller registers,
        memory-mapped I/O, or blocking operating system calls.
  \item \textbf{Tier 4: Application \& Telemetry Layer:} Front-panel gesture translators,
        graphical LCD rendering, and remote serial console handlers.
\end{enumerate}

\begin{figure}[H]
  \centering
  \resizebox{\linewidth}{!}{%
  \begin{tikzpicture}[
      m/.style={draw, thick, rounded corners, align=center, font=\small, minimum height=0.85cm, fill=white, inner sep=3pt},
      a/.style={-{Stealth[length=2.4mm]}, thick},
      layer/.style={draw=blue!60!black, dashed, rounded corners=4pt, fill=none},
      layer_hdr/.style={draw=blue!70!black, thick, rounded corners=3pt, fill=white, font=\footnotesize\bfseries\color{blue!80!black}, align=center, text width=2.1cm, inner sep=2pt},
    ]
    % Physical Inputs (top)
    \node[m, fill=white, text width=11.0cm] (in) at (1.4, 0.0) 
      {\textbf{Physical Inputs}: Onboard Key K1 (PC13), Ext Keys K0/STOP, Coin Pulses, UART Stream};

    % Layer 1: Input HAL Drivers (y = -1.9)
    \node[m, text width=3.1cm] (deb) at (-2.3, -1.9) 
      {\textbf{Key Debouncer}\\[2pt]\footnotesize $30\,\text{ms}$ Glitch Filter};
    \node[m, text width=3.1cm] (gst) at (1.4, -1.9) 
      {\textbf{Gesture Decoder}\\[2pt]\footnotesize 5-State Sub-FSM};
    \node[m, dashed, text width=3.0cm] (coin) at (5.1, -1.9) 
      {\textbf{Pulse Validator}\\[2pt]\footnotesize FIFO Coin Accumulator};

    % Layer 2: Application Front-Panel Dispatcher (y = -4.3)
    \node[m, text width=8.6cm] (evt) at (1.4, -4.3) 
      {\textbf{Front-Panel Event Dispatcher}\\[2pt]\footnotesize Translates Gestures \& Pulses to Formal Events (COIN, RUN, PAUSE, STOP)};

    % Layer 3: Core Logic & Prescaler Time-Base (y = -6.65)
    \node[m, text width=2.7cm] (tick) at (-2.9, -6.65) 
      {\textbf{Time-Base Engine}\\[2pt]\footnotesize SysTick ($1\,\text{ms} \to 1\,\text{s}$)};
    \node[m, fill=stRun, text width=5.4cm, font=\small\bfseries] (fsm) at (3.5, -6.65) 
      {Deterministic EFSM Core\\[2pt]\footnotesize\mdseries 6 Modes $\times$ Tuple $\mathbf{c} = (b, t, n, w, f)$};

    % Layer 4: Output HAL Callbacks & Drivers (y = -9.1 and y = -11.3)
    \node[m, text width=8.6cm] (cb) at (1.4, -9.1) 
      {\textbf{HAL Output Callback Interface}\\[2pt]\footnotesize Strictly Typed Actuator \& Optical Callback Table};
    \node[m, text width=3.1cm] (led) at (-2.3, -11.3) 
      {\textbf{LED Blinker Engine}\\[2pt]\footnotesize $1\,\text{Hz} / 2\,\text{Hz}$ Square-Wave};
    \node[m, text width=3.1cm] (act) at (1.4, -11.3) 
      {\textbf{Actuator Controller}\\[2pt]\footnotesize Motor Tumbler / Drain Pump};
    \node[m, text width=3.0cm] (ui) at (5.1, -11.3) 
      {\textbf{Dashboard \& Console}\\[2pt]\footnotesize ST7735 SPI4 \& USART1};

    % Physical Board Peripherals (bottom)
    \node[m, fill=white, text width=11.0cm] (hw) at (1.4, -13.5) 
      {\textbf{Target Microcontroller Hardware}: GPIO Registers, SPI4 Master, USART1 Telemetry, TIM1 PWM, SysTick};

    % Architectural layer grouping background boxes
    \begin{scope}[on background layer]
      \draw[layer] (-6.9, -1.0) rectangle (6.9, -2.8);
      \draw[layer] (-6.9, -3.3) rectangle (6.9, -5.3);
      \draw[layer] (-6.9, -5.5) rectangle (6.9, -7.8);
      \draw[layer] (-6.9, -8.2) rectangle (6.9, -12.4);
    \end{scope}

    % Dedicated Left Column Layer Headers (never intersects vertical signal arrows)
    \node[layer_hdr] at (-5.6, -1.9) {\textbf{Layer 1}\\[2pt]\footnotesize Input Drivers};
    \node[layer_hdr] at (-5.6, -4.3) {\textbf{Layer 2}\\[2pt]\footnotesize Dispatcher};
    \node[layer_hdr] at (-5.6, -6.65) {\textbf{Layer 3}\\[2pt]\footnotesize EFSM Core};
    \node[layer_hdr] at (-5.6, -10.3) {\textbf{Layer 4}\\[2pt]\footnotesize Output HAL\\[1pt]\footnotesize Callbacks};

    % Signal flows
    \draw[a] (in.south -| deb.north) -- (deb.north);
    \draw[a] (in.south -| coin.north) -- (coin.north);
    \draw[a] (deb) -- (gst);
    \draw[a] (gst) -- (gst |- evt.north);
    \draw[a, dashed] (coin.south) -- (coin.south |- evt.north);
    \draw[a] (evt.south -| fsm.north) -- node[right=3pt, font=\footnotesize, fill=white, inner sep=2pt, pos=0.75]{Events} (fsm.north);
    \draw[a] (tick) -- node[above=2pt, font=\footnotesize, fill=white, inner sep=2pt]{$1\,\text{s}$ Ticks} (fsm);
    \draw[a] (fsm.south) -- (fsm.south |- cb.north);
    \draw[a] (cb.south -| led.north) -- (led.north);
    \draw[a] (cb.south -| act.north) -- (act.north);
    \draw[a] (cb.south -| ui.north) -- (ui.north);
    \draw[a] (led.south) -- (led.south |- hw.north);
    \draw[a] (act.south) -- (act.south |- hw.north);
    \draw[a] (ui.south) -- (ui.south |- hw.north);
  \end{tikzpicture}%
  }
  \caption{Layered software architecture of the washing machine control system}
  \label{fig:arch}
\end{figure}

\subsection{Modular Decomposition and Source File Responsibilities}

Table~\ref{tab:modules} catalogues the responsibilities, dependencies, and interfaces of every
source module comprising the firmware codebase.

\begin{longtable}{L{4.8cm}L{10.8cm}}
  \caption{Software modules and component breakdown}\label{tab:modules}\\
  \toprule
  \textbf{Module (File)} & \textbf{Architectural Responsibility \& Implementation Details} \\
  \midrule
  \endfirsthead
  \toprule
  \textbf{Module (File)} & \textbf{Architectural Responsibility \& Implementation Details} \\
  \midrule
  \endhead
  \textbf{Low-Level Board Drivers}\newline\code{board_h750.c} &
    Configures microcontroller clock tree (400\,MHz), GPIO alternate functions, SysTick reload register,
    and power domain VOS1. Provides raw pin read functions. Contains zero application logic. \\
  \textbf{Button Debounce Engine}\newline\code{hal_button_engine.c} &
    Implements a symmetric $30\,\text{ms}$ software integrator for each physical key. Raw pin state
    must remain stable for 30 consecutive 1\,ms ticks before latching an edge, eliminating mechanical chatter. \\
  \textbf{Multi-Gesture Decoder}\newline\code{hal_button_gesture.c} &
    Internal 5-state sub-FSM (\code{IDLE}, \code{DOWN1}, \code{WAIT}, \code{DOWN2}, \code{HELD}) converting
    debounced pulses into high-level gesture tokens: Single Click, Double Click, Hold 0.6\,s, and Hold 1.2\,s. \\
  \textbf{Pulse Coin Validator}\newline\code{hal_coin_pulse.c} &
    Asynchronous pulse-train decoder for industrial coin acceptors. Counts pulse bursts ($1 = 10\,\text{\textcent}$,
    $2 = 20\,\text{\textcent}$, $5 = 50\,\text{\textcent}$) with $150\,\text{ms}$ inter-coin timeout detection and a 4-coin FIFO. \\
  \textbf{Front Panel Manager}\newline\code{wm_two_button.c} &
    Translates user gestures into formal FSM events conditioned upon the current qualitative state. Manages
    deposit staging and cancellation feedback. \\
  \textbf{FSM Controller Core}\newline\code{washing_machine_fsm.c} &
    Encapsulates the formal 6-state EFSM. Implements single-entry dispatcher \code{wm_fsm_dispatch_event()},
    guarded transitions, and context mutations. 100\% decoupled from target hardware registers. \\
  \textbf{Timer \& Clock Scheduler}\newline\code{main_h750.c} &
    Processes the SysTick 1\,ms interrupt counter in a non-blocking main loop catch-up algorithm, synthesizing
    jitter-free $1\,\text{ms}$ and $1\,\text{s}$ time bases without drift. \\
  \textbf{LED Blinker Engine}\newline\code{hal_led_blinker.c} &
    Non-blocking square-wave waveform generator: solid OFF, solid ON, $1.0\,\text{Hz}$ ($500\,\text{ms}$ ON / $500\,\text{ms}$ OFF),
    and $2.0\,\text{Hz}$ ($250\,\text{ms}$ ON / $250\,\text{ms}$ OFF). Enforces actuator safety invariant checks. \\
  \textbf{Display Driver \& Telemetry}\newline\code{lcd_st7735.c, gfx.c, ui_render.c} &
    SPI4 driver, 2D graphic primitive pipeline, and full-color 160$\times$80 dashboard renderer. Operates on a
    RAM back-buffer with event-driven dirty-flag optimization. \\
  \textbf{Serial CLI Console}\newline\code{ui_text.c} &
    Zero-heap telemetry formatter streaming human-readable ASCII status lines over $115\,200\,\text{bps}$ USART1
    and decoding single-key PC keyboard control commands. \\
  \bottomrule
\end{longtable}

\subsection{Decoupling via the Dependency Inversion Principle}

The FSM communicates with the physical environment exclusively through a strongly typed callback table
(\code{hal_output_callbacks_t}), passed as a pointer during \code{wm_fsm_init()}:

\begin{lstlisting}[caption={Formal HAL output callback interface definition}]
typedef struct {
    void (*set_rled)(hal_led_state_t state);
    void (*set_bled)(hal_led_state_t state);
    void (*set_motor)(hal_motor_state_t state);
    void (*set_water_valve)(bool open);
    void (*set_drain_pump)(bool on);
    void (*set_door_lock)(bool locked);
    void (*on_cycle_complete)(void);
    void (*return_coins)(uint32_t cents);   /* optional pre-run refund */
} hal_output_callbacks_t;
\end{lstlisting}

\noindent This structural decoupling provides three decisive advantages:
\begin{enumerate}
  \item \textbf{100\% Host Testability:} The exact same compiled FSM core object code
        is linked into the automated test runner on Linux/Windows (`gcc`), verifying all 34 transition rules
        in $<0.15\,\text{s}$ with zero hardware dependencies.
  \item \textbf{Microcontroller Agnosticism:} The FSM core can be ported to any target MCU (e.g., STM32F4,
        ESP32, AVR, or RISC-V) simply by supplying a platform-specific callback table.
  \item \textbf{Defensive Null-Pointer Guarding:} In compliance with MISRA-C Rule 1.3, every callback invocation
        is strictly guarded against NULL pointers before execution:
        \begin{center}
          \small\code{if ((fsm->callbacks != NULL) && (fsm->callbacks->set_motor != NULL)) { fsm->callbacks->set_motor(cmd); }}
        \end{center}
\end{enumerate}

\subsection{Event Dispatch Sequence}

Figure~\ref{fig:seq-diagram} depicts the end-to-end execution sequence when a user depresses the physical
\code{K1} push-button to start a washing cycle from \st{READY}.

\begin{figure}[H]
  \centering
  \resizebox{\linewidth}{!}{%
  \begin{tikzpicture}[
      font=\small,
      lifeline/.style={draw=gray!50, thick, dashed},
      box/.style={draw, thick, rounded corners, fill=white, minimum width=1.9cm, text width=1.9cm, minimum height=0.8cm, align=center, font=\footnotesize\bfseries},
      call/.style={thick, -{Stealth[length=2mm]}},
      ret/.style={thick, dashed, -{Stealth[length=2mm]}},
      note/.style={draw=orange!80!black, thick, fill=white, rounded corners, font=\footnotesize, align=center, inner sep=4pt}
    ]
    % Entities (lifelines spaced evenly across 14.0cm textwidth)
    \node[box] (usr) at (0.5,0) {User /\\Plant};
    \node[box] (gpio) at (3.0,0) {GPIO Driver\\PC13};
    \node[box] (deb) at (5.5,0) {Debounce\\Engine};
    \node[box] (gst) at (8.0,0) {Gesture\\Decoder};
    \node[box, fill=stRun] (fsm) at (10.8,0) {FSM Core\\EFSM};
    \node[box, fill=teal!10] (hal) at (13.8,0) {HAL Callbacks\\Outputs};

    % Lifelines
    \draw[lifeline] (usr) -- (0.5,-9.6);
    \draw[lifeline] (gpio) -- (3.0,-9.6);
    \draw[lifeline] (deb) -- (5.5,-9.6);
    \draw[lifeline] (gst) -- (8.0,-9.6);
    \draw[lifeline] (fsm) -- (10.8,-9.6);
    \draw[lifeline] (hal) -- (13.8,-9.6);

    % Call invocations (ample clearance below entity headers)
    \draw[call] (0.5,-1.4) -- node[above=2pt, font=\footnotesize, fill=white, inner sep=2pt]{Press K1} (3.0,-1.4);
    \draw[call] (3.0,-2.0) -- node[above=2pt, font=\footnotesize, fill=white, inner sep=2pt]{Raw HIGH} (5.5,-2.0);
    \draw[call] (5.5,-2.8) -- node[above=2pt, font=\footnotesize, fill=white, inner sep=2pt]{Stable Press ($30\,\text{ms}$)} (8.0,-2.8);
    \draw[call] (0.5,-3.4) -- node[above=2pt, font=\footnotesize, fill=white, inner sep=2pt]{Release K1} (3.0,-3.4);
    \draw[call] (8.0,-4.2) -- node[above=2pt, font=\footnotesize, fill=white, inner sep=2pt]{Single Click Token} (10.8,-4.2);
    
    % FSM Transition Execution Note
    \node[note, text width=3.4cm] at (10.8,-5.3) 
      {\textbf{T08 Fired}: \st{READY} $\to$ \st{RUNNING}\\[2pt]$b \gets 0, \quad t \gets 1800\,\text{s}$};

    % Hardware-abstracted actuator commands (clean labels, generous travel)
    \draw[call] (10.8,-6.7) -- node[above=2pt, font=\footnotesize, fill=white, inner sep=2pt]{Lock Door Deadbolt} (13.8,-6.7);
    \draw[call] (10.8,-7.6) -- node[above=2pt, font=\footnotesize, fill=white, inner sep=2pt]{Command Drum Agitate (PWM)} (13.8,-7.6);
    \draw[call] (10.8,-8.5) -- node[above=2pt, font=\footnotesize, fill=white, inner sep=2pt]{Command BLED ($1.0\,\text{Hz}$ Blink)} (13.8,-8.5);
  \end{tikzpicture}%
  }
  \caption{UML sequence diagram: Deterministic signal propagation from physical press to actuator callback}
  \label{fig:seq-diagram}
\end{figure}
